% Página 37 do PDF de Gabriel
\begin{frame}[plain,noframenumbering,label=original-37]
  \begin{tikzpicture}[remember picture,overlay,x=1bp,y=-1bp]
  \begin{scope}[shift={(current page.north west)}]
  \Fundo{37}
  \Texto{44.456}{41.894}{90.657}{8.518}{b}{34E0C8}{ANÁLISE DESCRITIVA}
  \Texto{43.039}{116.125}{63.147}{56.787}{b}{FFB014}{06}
  \Texto{43.039}{154.533}{202.078}{24.608}{b}{FFFFFF}{Correlação linear}
  \Texto{43.039}{234.662}{46.285}{7.572}{b}{34E0C8}{EMASC 2026}
  \Texto{91.429}{234.662}{2.65}{7.572}{b}{FFB014}{\textbullet{}}
  \Texto{96.176}{234.662}{21.45}{7.572}{b}{34E0C8}{LAME}
  \Link{42.042}{135.586}{204.068}{25.418}{original-37}
  \end{scope}
  \end{tikzpicture}
\end{frame}


% Página 38 do PDF de Gabriel
\begin{frame}[plain,noframenumbering,label=original-38]
  \begin{tikzpicture}[remember picture,overlay,x=1bp,y=-1bp]
  \begin{scope}[shift={(current page.north west)}]
  \Fundo{38}
  \Texto{40.175}{27.803}{236.128}{17.983}{b}{2C2084}{Por que medir associação?}
  \Texto{35.433}{63.506}{307.909}{8.518}{i}{737696}{Duas variáveis podem variar juntas; queremos resumir isso em um único número.}
  \Texto{38.534}{100.643}{34.716}{9.465}{b}{027499}{Direção}
  \Texto{38.534}{115.586}{163.708}{9.465}{r}{000000}{Quando uma variável aumenta, a outra}
  \Texto{38.534}{127.542}{136.156}{9.465}{r}{000000}{tende a aumentar ou a diminuir?}
  \Texto{250.943}{100.643}{52.594}{9.465}{b}{CC8D10}{Intensidade}
  \Texto{250.943}{115.587}{143.084}{9.465}{r}{000000}{Quão próxima de uma reta é essa}
  \Texto{250.943}{127.542}{36.296}{9.465}{r}{000000}{relação?}
  \Texto{49.227}{162.195}{153.528}{8.518}{r}{000000}{O coeficiente de correlação de Pearson,}
  \Texto{205.623}{162.195}{2.836}{8.518}{i}{000000}{r}
  \Texto{209.418}{162.195}{196.894}{8.518}{r}{000000}{, resume esses dois aspectos em um número entre}
  \Texto{49.227}{184.113}{55.87}{8.518}{r}{000000}{saúde mental?}
  \Texto{37.21}{247.463}{40.499}{6.625}{b}{737696}{EMASC 2026}
  \Texto{80.653}{247.463}{2.319}{6.625}{r}{FFB014}{\textbullet{}}
  \Texto{85.913}{247.463}{117.821}{6.625}{r}{737696}{Liga Acadêmica Matemática \& Empresa}
  \Texto{206.682}{247.463}{2.319}{6.625}{r}{FFB014}{\textbullet{}}
  \Texto{211.943}{247.463}{13.25}{6.625}{r}{737696}{UFS}
  \Texto{397.549}{248.132}{7.367}{6.625}{b}{2C2084}{24}
  \Texto{406.02}{248.132}{10.315}{6.625}{r}{737696}{/ 40}
  \end{scope}
  \end{tikzpicture}
\end{frame}


% Página 39 do PDF de Gabriel
\begin{frame}[plain,noframenumbering,label=original-39]
  \begin{tikzpicture}[remember picture,overlay,x=1bp,y=-1bp]
  \begin{scope}[shift={(current page.north west)}]
  \Fundo{39}
  \Texto{40.175}{27.803}{197.052}{17.983}{b}{2C2084}{Correlação de Pearson}
  \Texto{35.433}{63.506}{225.523}{8.518}{i}{737696}{A fórmula traduz a associação em um número padronizado.}
  \Texto{46.238}{128.971}{107.397}{6.974}{t}{000000}{a\ =\ np.array([1,\ 2,\ 3,\ 4,\ 5])}
  \Texto{46.238}{137.938}{107.397}{6.974}{t}{000000}{b\ =\ np.array([2,\ 1,\ 4,\ 3,\ 5])}
  \Texto{46.238}{146.904}{129.622}{6.974}{t}{000000}{da,\ db\ =\ a\ -\ a.mean(),\ b\ -\ b.mean()}
  \Texto{46.238}{155.87}{51.843}{6.974}{t}{000000}{num\ =\ (da*db).}
  \Texto{98.081}{155.87}{11.116}{6.974}{t}{2C2084}{sum}
  \Texto{109.197}{155.87}{7.406}{6.974}{t}{000000}{()}
  \Texto{212.915}{155.87}{11.109}{6.974}{c}{737696}{\#\ 8}
  \Texto{46.238}{164.837}{81.475}{6.974}{t}{000000}{den\ =\ np.sqrt((da**2).}
  \Texto{127.713}{164.837}{11.109}{6.974}{t}{2C2084}{sum}
  \Texto{138.822}{164.837}{40.734}{6.974}{t}{000000}{()*(db**2).}
  \Texto{179.556}{164.837}{11.109}{6.974}{t}{2C2084}{sum}
  \Texto{190.665}{164.837}{11.109}{6.974}{t}{000000}{())}
  \Texto{212.915}{164.837}{14.812}{6.974}{c}{737696}{\#\ 10}
  \Texto{46.238}{173.803}{33.328}{6.974}{t}{000000}{num\ /\ den}
  \Texto{212.915}{173.803}{18.515}{6.974}{c}{737696}{\#\ 0.8}
  \Texto{46.238}{182.769}{85.178}{6.974}{t}{000000}{np.corrcoef(a,\ b)[0,\ 1]}
  \Texto{212.915}{182.769}{18.515}{6.974}{c}{737696}{\#\ 0.8}
  \Texto{49.227}{217.319}{350.78}{8.518}{r}{000000}{O numerador soma os produtos dos desvios; o denominador padroniza o resultado. Por isso}
  \Texto{402.515}{217.319}{2.836}{8.518}{i}{000000}{r}
  \Texto{37.21}{247.463}{40.499}{6.625}{b}{737696}{EMASC 2026}
  \Texto{80.653}{247.463}{2.319}{6.625}{r}{FFB014}{\textbullet{}}
  \Texto{85.913}{247.463}{117.821}{6.625}{r}{737696}{Liga Acadêmica Matemática \& Empresa}
  \Texto{206.682}{247.463}{2.319}{6.625}{r}{FFB014}{\textbullet{}}
  \Texto{211.943}{247.463}{13.25}{6.625}{r}{737696}{UFS}
  \Texto{397.549}{248.132}{7.367}{6.625}{b}{2C2084}{25}
  \Texto{406.02}{248.132}{10.315}{6.625}{r}{737696}{/ 40}
  \end{scope}
  \end{tikzpicture}
\end{frame}


% Páginas 40 a 43 do PDF de Gabriel
\begin{frame}[plain,noframenumbering,label=original-40]
  \begin{tikzpicture}[remember picture,overlay,x=1bp,y=-1bp]
  \begin{scope}[shift={(current page.north west)}]
  \only<1>{\Fundo{40}}
  \only<2>{\Fundo{41}}
  \only<3>{\Fundo{42}}
  \only<4>{\Fundo{43}}
  \Texto{40.175}{27.803}{221.832}{17.983}{b}{2C2084}{Limitações do coeficiente}
  \Texto{35.433}{63.506}{240.361}{8.518}{i}{737696}{r resume associação linear; fora desse contexto, pode enganar.}
  \only<1>{\Texto{49.227}{207.561}{22.036}{8.518}{r}{000000}{sável.}}
  \Texto{37.21}{247.463}{40.499}{6.625}{b}{737696}{EMASC 2026}
  \Texto{80.653}{247.463}{2.319}{6.625}{r}{FFB014}{\textbullet{}}
  \Texto{85.913}{247.463}{117.821}{6.625}{r}{737696}{Liga Acadêmica Matemática \& Empresa}
  \Texto{206.682}{247.463}{2.319}{6.625}{r}{FFB014}{\textbullet{}}
  \Texto{211.943}{247.463}{13.25}{6.625}{r}{737696}{UFS}
  \Texto{397.549}{248.132}{7.367}{6.625}{b}{2C2084}{26}
  \Texto{406.02}{248.132}{10.315}{6.625}{r}{737696}{/ 40}
  \only<2,3>{\Texto{49.227}{196.603}{7.104}{8.518}{b}{000000}{2.}}
  \only<2>{\Texto{59.261}{196.603}{236.443}{8.518}{r}{000000}{Sem o ponto destacado, não há associação entre as variáveis.}}
  \only<3>{\Texto{59.261}{196.603}{191.553}{8.518}{r}{000000}{Um único ponto extremo cria uma correlação forte.}}
  \only<4>{\Texto{49.227}{196.603}{7.104}{8.518}{b}{000000}{3.}}
  \only<4>{\Texto{59.185}{196.603}{347.126}{8.518}{r}{000000}{O calor aumenta as vendas de sorvete e o número de afogamentos. A associação entre eles}}
  \only<4>{\Texto{49.227}{207.561}{50.665}{8.518}{r}{000000}{não é causal.}}
  \end{scope}
  \end{tikzpicture}
\end{frame}


% Página 44 do PDF de Gabriel
\begin{frame}[plain,noframenumbering,label=original-44]
  \begin{tikzpicture}[remember picture,overlay,x=1bp,y=-1bp]
  \begin{scope}[shift={(current page.north west)}]
  \Fundo{44}
  \Texto{40.175}{27.803}{190.866}{17.983}{b}{2C2084}{Correlação nos dados}
  \Texto{35.433}{63.506}{221.689}{8.518}{i}{737696}{Retomamos a pergunta inicial com o coeficiente calculado.}
  \Texto{49.227}{214.786}{357.083}{8.518}{r}{000000}{negativa de intensidade moderada. O coeficiente não indica quantos pontos se perdem por}
  \Texto{49.227}{225.745}{186.817}{8.518}{r}{000000}{hora de uso; mede apenas direção e intensidade.}
  \Texto{37.21}{247.463}{40.499}{6.625}{b}{737696}{EMASC 2026}
  \Texto{80.653}{247.463}{2.319}{6.625}{r}{FFB014}{\textbullet{}}
  \Texto{85.913}{247.463}{117.821}{6.625}{r}{737696}{Liga Acadêmica Matemática \& Empresa}
  \Texto{206.682}{247.463}{2.319}{6.625}{r}{FFB014}{\textbullet{}}
  \Texto{211.943}{247.463}{13.25}{6.625}{r}{737696}{UFS}
  \Texto{397.549}{248.132}{7.367}{6.625}{b}{2C2084}{27}
  \Texto{406.02}{248.132}{10.315}{6.625}{r}{737696}{/ 40}
  \end{scope}
  \end{tikzpicture}
\end{frame}


% Páginas 45 a 47 do PDF de Gabriel
\begin{frame}[plain,noframenumbering,label=original-45]
  \begin{tikzpicture}[remember picture,overlay,x=1bp,y=-1bp]
  \begin{scope}[shift={(current page.north west)}]
  \only<1>{\Fundo{45}}
  \only<2>{\Fundo{46}}
  \only<3>{\Fundo{47}}
  \Texto{40.175}{27.803}{185.867}{17.983}{b}{2C2084}{Matriz de correlações}
  \Texto{35.433}{63.506}{281.179}{8.518}{i}{737696}{O mesmo coeficiente pode ser calculado para todos os pares de variáveis.}
  \Texto{254.928}{93.035}{17.982}{8.518}{b}{132482}{Azul}
  \Texto{272.909}{93.035}{2.368}{8.518}{r}{000000}{:}
  \Texto{282.143}{93.035}{81.287}{8.518}{r}{000000}{associação negativa.}
  \Texto{372.259}{93.035}{38.391}{8.518}{b}{E27800}{Vermelho}
  \Texto{410.649}{93.035}{2.368}{8.518}{r}{000000}{:}
  \Texto{254.928}{103.994}{76.347}{8.518}{r}{000000}{associação positiva.}
  \Texto{335.224}{103.994}{29.345}{8.518}{b}{000000}{Branco}
  \Texto{364.569}{103.994}{48.45}{8.518}{r}{000000}{: próximo de}
  \Texto{254.928}{114.953}{18.467}{8.518}{r}{000000}{zero.}
  \Texto{37.209}{247.463}{40.499}{6.625}{b}{737696}{EMASC 2026}
  \Texto{80.653}{247.463}{2.319}{6.625}{r}{FFB014}{\textbullet{}}
  \Texto{85.914}{247.463}{117.821}{6.625}{r}{737696}{Liga Acadêmica Matemática \& Empresa}
  \Texto{206.683}{247.463}{2.319}{6.625}{r}{FFB014}{\textbullet{}}
  \Texto{211.943}{247.463}{13.25}{6.625}{r}{737696}{UFS}
  \Texto{397.549}{248.132}{7.367}{6.625}{b}{2C2084}{28}
  \Texto{406.02}{248.132}{10.315}{6.625}{r}{737696}{/ 40}
  \only<2,3>{\Texto{254.928}{140.933}{158.086}{8.518}{r}{000000}{A diagonal vale 1, pois cada variável é per-}}
  \only<2,3>{\Texto{254.928}{151.892}{158.086}{8.518}{r}{000000}{feitamente associada a si mesma, e a ma-}}
  \only<2,3>{\Texto{254.928}{162.851}{57.522}{8.518}{r}{000000}{triz é simétrica.}}
  \only<3>{\Texto{254.928}{188.635}{158.086}{8.518}{r}{000000}{A idade não apresenta associação com as}}
  \only<3>{\Texto{254.928}{199.594}{158.086}{8.518}{r}{000000}{demais variáveis, o que reforça a hipótese}}
  \only<3>{\Texto{254.928}{210.553}{77.906}{8.518}{r}{000000}{de dados simulados.}}
  \end{scope}
  \end{tikzpicture}
\end{frame}


% Páginas 48 a 51 do PDF de Gabriel
\begin{frame}[plain,noframenumbering,label=original-48]
  \begin{tikzpicture}[remember picture,overlay,x=1bp,y=-1bp]
  \begin{scope}[shift={(current page.north west)}]
  \only<1>{\Fundo{48}}
  \only<2>{\Fundo{49}}
  \only<3>{\Fundo{50}}
  \only<4>{\Fundo{51}}
  \Texto{40.175}{27.803}{224.871}{17.983}{b}{2C2084}{Associação e causalidade}
  \Texto{35.433}{63.506}{203.12}{8.518}{i}{737696}{Medir associação não responde à pergunta de causa.}
  \Texto{71.408}{84.712}{22.726}{8.518}{b}{262848}{redes}
  \Texto{55.552}{127.429}{54.439}{8.518}{b}{262848}{saúde mental}
  \Texto{35.302}{147.661}{94.938}{9.465}{r}{262848}{Redes afetam a saúde}
  \Texto{65.939}{159.616}{33.665}{9.465}{r}{262848}{mental?}
  \Texto{37.21}{247.464}{40.499}{6.625}{b}{737696}{EMASC 2026}
  \Texto{80.653}{247.464}{2.319}{6.625}{r}{FFB014}{\textbullet{}}
  \Texto{85.913}{247.464}{117.821}{6.625}{r}{737696}{Liga Acadêmica Matemática \& Empresa}
  \Texto{206.682}{247.464}{2.319}{6.625}{r}{FFB014}{\textbullet{}}
  \Texto{211.942}{247.464}{13.25}{6.625}{r}{737696}{UFS}
  \Texto{397.549}{248.133}{7.367}{6.625}{b}{2C2084}{29}
  \Texto{406.02}{248.133}{10.315}{6.625}{r}{737696}{/ 40}
  \only<2,3,4>{\Texto{215.409}{84.713}{22.726}{8.518}{b}{262848}{redes}}
  \only<2,3,4>{\Texto{199.553}{127.429}{54.439}{8.518}{b}{262848}{saúde mental}}
  \only<2,3,4>{\Texto{167.899}{147.661}{117.748}{9.465}{r}{262848}{A saúde mental muda o uso}}
  \only<2,3,4>{\Texto{205.728}{159.616}{42.089}{9.465}{r}{262848}{de redes?}}
  \only<3,4>{\Texto{355.394}{84.896}{26.977}{8.518}{b}{262848}{fator X}}
  \only<3,4>{\Texto{332.008}{132.993}{22.726}{8.518}{b}{262848}{redes}}
  \only<3,4>{\Texto{384.92}{133.191}{24.617}{8.518}{b}{262848}{saúde}}
  \only<3,4>{\Texto{317.668}{153.422}{106.211}{9.465}{r}{262848}{Um terceiro fator afeta as}}
  \only<3,4>{\Texto{357.882}{165.377}{25.781}{9.465}{r}{262848}{duas?}}
  \only<4>{\Texto{49.227}{185.544}{118.289}{8.518}{r}{000000}{Os três diagramas produzem o}}
  \only<4>{\Texto{170.3}{185.544}{29.822}{8.518}{b}{000000}{mesmo}}
  \only<4>{\Texto{202.904}{185.544}{203.402}{8.518}{r}{000000}{coeficiente de correlação. Dados observacionais, por}}
  \only<4>{\Texto{49.227}{196.503}{152.992}{8.518}{r}{000000}{si só, não permitem distinguir entre eles.}}
  \end{scope}
  \end{tikzpicture}
\end{frame}
